%======================================================================
\chapter{Introduction and plan of the notes}\label{ch:intro}
%======================================================================

\section{The theorem these notes are about}

A \emph{continuum} is a nonempty compact connected metric space. Continua include the
closed interval $[0,1]$, the circle, the closed disc, and many stranger objects. In this
course a continuum will usually sit inside the plane $\R^2$, and we call it a \emph{plane
continuum}. A continuum is \emph{nondegenerate} if it has more than one point.

A space $M$ is \emph{homogeneous} if for any two points $x,y\in M$ there is a
homeomorphism $h\colon M\to M$ with $h(x)=y$. In words: every point looks like every
other point. The circle is homogeneous (rotate it). The interval $[0,1]$ is not: removing
an end point leaves a connected set, while removing an interior point does not, and a
homeomorphism cannot change this.

The final goal of the whole series of notes is the following theorem, which completes a
programme begun in the 1920s.

\begin{theorem}[Classification of homogeneous plane continua]\label{thm:classification}
A nondegenerate plane continuum is homogeneous if and only if it is homeomorphic to one
of the following three continua:
\begin{enumerate}[label=(\roman*)]
\item the circle;
\item the pseudo-arc;
\item the circle of pseudo-arcs.
\end{enumerate}
Each of the three is realized in the plane, each is homogeneous, and no two of them are
homeomorphic.
\end{theorem}

The theorem has been proved in full in the Lean~4 proof assistant, in the repository
that accompanies these notes. A proof assistant checks every logical step of a proof
down to the axioms. The formal proof is therefore reliable, but it is not written to be
read by people: it consists of roughly a thousand files of code. These notes translate
it back into ordinary mathematics. They follow the formal proof step by step, and where
the formal proof departs from the original papers (it often does, to repair a gap or to
find a route that is easier to check) the notes follow the formal proof.

\section{The three continua, informally}

\paragraph{The circle.} Nothing needs to be said, except that ``circle'' means any space
homeomorphic to $S^1=\set{z\in\mathbb{C} : \abs{z}=1}$; such a space is also called a
\emph{simple closed curve}.

\paragraph{The pseudo-arc.} Here is a first description. Cover a continuum $M$ by a finite
sequence of open sets $L_0,\dots,L_{n-1}$ (``links'') such that consecutive links meet and
non-consecutive links do not. Such a cover is called a \emph{chain}. A continuum that can
be covered by chains with arbitrarily small links is called \emph{snake-like} (or
\emph{chainable}). The interval $[0,1]$ is snake-like: cover it by short overlapping
intervals. The circle is not snake-like.

A continuum is \emph{decomposable} if it is the union of two proper subcontinua, and
\emph{indecomposable} otherwise. The interval is decomposable ($[0,1]=[0,\frac12]\cup
[\frac12,1]$). It is a surprising fact that indecomposable nondegenerate continua exist at
all. A continuum is \emph{hereditarily indecomposable} if every subcontinuum is
indecomposable.

\begin{definition}\label{def:pseudoarc-intro}
A \emph{pseudo-arc} is a nondegenerate continuum that is snake-like and hereditarily
indecomposable.
\end{definition}

That such a continuum exists is the main theorem of this volume. That any two pseudo-arcs
are homeomorphic, and that the pseudo-arc is homogeneous, are the main theorems of
Volume~2. The pseudo-arc is a very strange object: it contains no arc (no subset
homeomorphic to $[0,1]$), yet it is homeomorphic to each of its nondegenerate
subcontinua.

\paragraph{The circle of pseudo-arcs.} This is a continuum $X$ with a continuous open map
$f\colon X\to S^1$ onto the circle, every fibre $f^{-1}(c)$ of which is a pseudo-arc, and
which is itself ``circle-like''. Think of a circle in which every point has been replaced
by a pseudo-arc. It is constructed and studied in Volume~8.

\section{The shape of the whole proof}

The proof of Theorem~\ref{thm:classification} splits as follows.
\begin{enumerate}
\item \textbf{The pseudo-arc.} A pseudo-arc exists in the plane; any two pseudo-arcs are
homeomorphic; the pseudo-arc is homogeneous.
\item \textbf{The indecomposable case.} A homogeneous indecomposable plane continuum is a
pseudo-arc. One shows in turn that it is hereditarily indecomposable, that it has ``span
zero'', and that a hereditarily indecomposable continuum of span zero is snake-like.
\item \textbf{The decomposable case.} A homogeneous decomposable plane continuum is either
a circle or admits an open map onto the circle whose fibres are pseudo-arcs.
\item \textbf{Circles of pseudo-arcs.} Any two continua of the last kind are homeomorphic,
they are homogeneous, and one of them lies in the plane.
\item \textbf{Assembly.}
\end{enumerate}
All of this rests on two bodies of general theory: the topology of continua (chains,
decompositions, the Effros theorem on homogeneous spaces), and the topology of the plane
(separation theorems, the Jordan curve theorem, and R.~L.~Moore's characterization of the
plane).

\section{Plan of the series}

The formal proof is far too long for one course. The notes are therefore planned in
volumes, each of which can be read on its own once its prerequisites are known. The plan
is as follows; only Volume~1 is written at present.

\begin{center}
\begin{tabular}{@{}lp{0.72\textwidth}@{}}
\toprule
Volume & Content\\
\midrule
1 & Chains and crookedness; descriptions; described continua are hereditarily
indecomposable; Bing's crooked pattern; construction of a pseudo-arc in the plane.
\emph{(This volume.)}\\
2 & Composants and indecomposable continua; Bing's consolidation theorem and pattern
following; matching covers and homeomorphisms; uniqueness of the pseudo-arc; homogeneity
of the pseudo-arc.\\
3 & General theory of continua: boundary bumping, irreducibility, non-cut points, triods,
snake-like and circle-like continua, upper semicontinuous decompositions, the
homeomorphism group and the Effros theorem.\\
4 & Plane topology I: continuous logarithms, Borsuk's theorem, Janiszewski's theorem,
polygons and the polygonal Jordan curve theorem, enclosing polygons, tree-like plane
continua.\\
5 & Plane topology II: Moore's axioms characterize the plane; decompositions of the plane
into non-separating continua; the Jordan curve theorem.\\
6 & The indecomposable case: consequences of homogeneity, hereditary indecomposability,
span zero, snake-likeness.\\
7 & The decomposable case: aposyndesis, Jones's aposyndetic decomposition, homogeneous
aposyndetic plane continua are circles.\\
8 & Circles of pseudo-arcs; the classification theorem.\\
\bottomrule
\end{tabular}
\end{center}

\section{What this volume proves}

The main theorem of this volume is:

\begin{theorem}[Existence of the pseudo-arc]\label{thm:main-vol1}
There is a pseudo-arc in the plane $\R^2$.
\end{theorem}

The proof has two halves.
\begin{itemize}
\item \emph{Abstract half} (Chapters~\ref{ch:prelim}--\ref{ch:descriptions}). We define
crookedness of one chain in another and \emph{descriptions}: sequences of chains, each
crooked in the one before and with links tending to zero in size. We prove that the
intersection of the levels of a description, if the links stay in a compact set, is a
hereditarily indecomposable snake-like continuum. This is a theorem of R.~H.~Bing from
1948.
\item \emph{Concrete half} (Chapters~\ref{ch:patterns}--\ref{ch:plane}). We write down an
explicit integer sequence, Bing's pattern, prove that it is crooked, and use it together
with an explicit homeomorphism of the plane to construct a description in $\R^2$ whose
links stay in a fixed rectangle.
\end{itemize}

The concrete half is where the formal proof differs most from the literature: the
classical construction is described with pictures, while the formal proof needs exact
coordinates. These notes give the exact coordinates and check every inequality, so that
the construction can be verified line by line without pictures.

\section{Prerequisites and conventions}

The reader should know the basics of metric spaces: open and closed sets, convergence,
continuity, compactness and connectedness. Everything beyond this is proved in the notes;
Chapter~\ref{ch:prelim} recalls what is used, with proofs of the less standard facts.

\begin{itemize}
\item $\N=\set{0,1,2,\dots}$ includes $0$. For $n\in\N$ we write $[n]=\set{0,1,\dots,n-1}$.
Chains are indexed from $0$, as in the formal proof.
\item The plane $\R^2$ carries the \emph{maximum metric}
$d\big((a,b),(c,d)\big)=\max\big(\abs{a-c},\abs{b-d}\big)$. It defines the usual topology.
We use it because it is the metric of the formal proof, and because it makes diameters of
rectangles easy to compute: the rectangle $[a,a']\times[b,b']$ has diameter
$\max(a'-a,b'-b)$.
\item For integers $a,b$ we write $\fl{a/b}$ for the integer part of $a/b$ ($a\ge0$, $b>0$).
\end{itemize}

\paragraph{Formal counterparts.} After many results there is a line ``Formal counterpart:''
giving the name of the corresponding statement in the Lean development. These names are
there for readers who want to compare with the formal proof; nothing in the notes depends
on them, and they may be ignored. Appendix~\ref{app:formal} collects them in one table.

\paragraph{Exercises.} Each chapter ends with exercises. None of them is used later.
